\documentclass[a4paper,12pt]{article}
\usepackage[utf8]{inputenc}
\usepackage[T1]{fontenc}
\usepackage[ngerman]{babel}
\usepackage{amsmath}
\usepackage{amssymb}
\usepackage{enumitem}
\usepackage{geometry}
\geometry{a4paper, margin=2.5cm}
\usepackage{parskip}

\title{Einsendeaufgaben -- Aufgabe 6.3\\[0.5em]
\large Modul 61112: Lineare Algebra}
\author{}
\date{}

\begin{document}
\maketitle

\section*{Aufgabe 6.3}

\begin{enumerate}[label=\alph*)]
  \item Da $\mathbb{C}$ algebraisch abgeschlossen ist, zerfällt $\chi_A$ in
  Linearfaktoren. Nach dem Hauptsatz zur Jordan'schen Normalform Satz 6.3.8
  ist $A$ ähnlich zu einer Matrix $J \in M_{n,n}(\mathbb{C})$ in Jordan'scher
  Normalform. Da die Ähnlichkeitsklasse eine Äquivalenzklasse und ähnlich zu
  $2A$ ist, ist auch $2A$ ähnlich zu $J$. Demnach besitzen $A$ und $2A$ nach
  Satz 6.3.9 die gleichen Eigenwerte.

  Das bedeutet aber, dass es nur einen Eigenwert von $A$ und $2A$ gibt, welcher
  die $0$ ist. Denn wenn wir einen beliebigen Eigenwert $\lambda \in \mathbb{C}$
  von $A$ wählen, folgt:
  \[
    2Av = 2\lambda v \quad \text{für alle } v \in \mathrm{Eig}_A(\lambda).
  \]
  Demnach ist zu jedem Eigenwert $\lambda \in \mathbb{C}$ von $A$, $2\lambda$ ein
  Eigenwert von $2A$. Da aber $A$ und $2A$ die gleichen Eigenwerte besitzen, muss
  also $0$ der einzige Eigenwert von $A$ und $2A$ sein.

  Da $0$ also der einzige Eigenwert von $A$ ist, sind die Jordan-Blöcke in
  nilpotenter Normalform und dadurch auch $J$. Abschließend können wir
  schlussfolgern, dass $A$ nilpotent ist.
\end{enumerate}

\end{document}
